\chapter{Objetos decidibles en Topos}


\begin{proposition}
  Los subespacios de un espacio discreto son discretos
\end{proposition}
\begin{proof}
  Sea $A$ un espacio discreto y $m\colon S\to A$ un subespacio de $A$, queremos demostrar que la 
  diagonal $\Delta_S\colon S\to S\times S$ es un subobjeto complementado. Por hipótesis, para $\Delta_A\colon A\to A\times A$
  existe su complemento, $C$, tal que $C\cap A=0$ y $C\cup A=A\times A$. Notamos que
  $m$ induce $m\times m\colon S\times S \rightarrowtail A\times A$ y $(m\times m)^{-1}(\Delta_A)=\Delta_S$ ya que el siguiente es un producto fibrado:
  \begin{equation*}
    \begin{tikzcd}
        %T\ar[bend left=30]{rrd}\ar[bend right=30]{rdd}[swap]{t}
          %\ar[dashed]{rd}{\pi_S t}\\[-2ex]
        S\ar{r}{m} \ar[tail]{d}[swap]{\Delta_S} & A\ar[tail]{d}{\Delta_A}\\
        S\times S\ar{r}[swap]{m\times m} & A\times A.
    \end{tikzcd}
  \end{equation*}
  Para demostrar que el diagrama es un producto fibrado necesitamos que el cuadrado conmute
  y que para cualquier objeto $T$ con flechas que haga conmutar al diagrama exterior, exista 
  un único morfismo $T\to S$ que haga conmutar los triángulos. Como $\Delta_A m = \langle \id_A, \id_A\rangle m = m\times m$ la conmutatividad
  del cuadrado es inmediata. Para un $T$ con flechas $t\colon T\to S\times S$, $u\colon T\to A$ tales
  que $\Delta_A u = \langle m, m\rangle t$, proponemos a $\pi_S t$ como la flecha de $T$
  a $S$, aprovechando que $t=\langle\pi_S t, \pi_S t\rangle$ donde $\pi_S\colon S\times S \to S$
  es la proyección a $S$. De esta manera tenemos la conmutatividad del triángulo izquierdo
  y, junto con la conmutatividad del cuadrado exterior e interior, obtenemos que $\Delta_A u
  = \Delta_A m \pi_S t$, pero $\Delta_A$ es mono, entonces $u=m \pi_S t$ y el triángulo superior
  también conmuta. La unicidad de $\pi_S t$ es inmediata ya que $\Delta_S$ es
  mono.\pdfcomment{Tal vez se pueda quitar la demostración}
  

  Dado que $S\times S$ es un subobjeto de $A\times A$ a través de $m\times m$, la imagen
  inversa $(m\times m)^{-1}$ es un morfismo de álgebras de Heyting $Sub(A\times A)\to Sub(S\times S)$,
  con lo cual $(m\times m)^{-1}$ preserva complementos y por lo anterior $\Delta_S\colon S\to S\times S$ es un
  subobjeto complementado con complemento $(m\times m)^{-1}(C)$.

\end{proof}

\begin{proposition}[label=prop:terminaldiscreto]
  El objeto terminal de un topos es un objeto decidible.
\end{proposition}
\begin{proof}
  Queremos mostrar que $\Delta_1\colon 1\to 1\times 1$ es un subobjeto complementado, pero $1\times 1 \cong 1$
  entonces $1$ es el subobjeto maximal, por lo que $0 \to 1\times 1$ es el complemento que buscamos.
\end{proof}

\begin{theorem}[label=thm:onndecidible]
  Sea $\mathcal{E}$ un topos con $(N,o,s)$ un objeto de números naturales, entonces $N$ es un objeto decidible.
\end{theorem}
\begin{proof}
  El plan es ver a $\Delta\colon N\to N\times N$ como intersección de subobjetos complementados en $\sub(N\times N)$ (así va la demostración de Johnstone pero hay una más directa con la lógica interna).
  Primero, en un topos tenemos recursividad con parámetros que nos permite definir a la resta truncada como 
  $\dot{-}\colon N\times N\to N$ tal que $\dot{-}(a,0)=a$ y $\dot{-}(a,n+1)=p(\dot{-}(a,n))$ donde $p$ es el predecesor de~\ref{def:predecesor}.
  Ahora construimos al subobjeto:
  \begin{equation*}
    \begin{tikzcd}
      N_1\ar{r}{}\ar[tail]{d}{} & 1\ar{d}{o} \\
      N\times N\ar{r}{\dot{-}} & N,
    \end{tikzcd}
  \end{equation*}
  que es complementado gracias a que ser complementado es estable bajo productos fibrados. Como $n\dot{-}n=0$ para toda $n$,
  $\Delta$ es subobjeto de $N_1$ y por lo tanto $\Delta$ también tiene complemento en $\sub(N\times N)$.
\end{proof}

\begin{proposition}
  Si $\toposE$ cumple el Nullstellensatz y $B$ es un objeto decidible, entonces para $f,g\colon A\to B$
  con $f\neq g$ existe un elemento global $p\colon 1\to A$ tal que $f p \neq g p$.
\end{proposition}
\begin{proof}
  Sea $\toposE$ un topos que satisface el Nullstellensatz, entonces $toposE$ es bivaluado y los soportes se escinden.
  Sea $B$ un objeto decidible y $f,g\colon A\to B$ con $f\neq g$. Consideramos al igualador $e\colon E\to A$ de $f$ y $g$, 
  quien de hecho es el producto fibrado entre $\delta\colon B\to B\times B$ y $\langle f,g\rangle\colon A\to B\times B$.
  Dado que $B$ es complementado en $\sub(B\times B)$, $E$ es complementado en $\sub(A)$, llamemos $\overline{e}\colon \overline{E}\to A$ a su complemento.
  Como $f\neq g$, $E\not\cong A$ y por lo tanto $\overline{E}\not\cong 0$. Por el Nullstellensatz, $\overline{E}$ tiene un elemento global $p\colon 1\to \overline{E}$ y por lo tanto $e p\colon 1\to A$ es un elemento global de $A$ tal que $f e p \neq g e p$.
\end{proof}







